\section{Roma and Amor}

\begin{quotex}
\textit{Omnia vincit amor et nos cedamus amori}
\flright{\textsc{Virgil}, \textit{Eclogue X}, 69}
\end{quotex}


The Roman poet Virgil composed this line: ``\textbf{Love conquers all and let us yield to love.}'' Vladimir Solovyov, changing the context somewhat, explains its significance:

\begin{quotationx}
Against the false man-god of political monarchy the true God-Man set up the spiritual power of ecclesiastical monarchy founded on truth and love. Universal monarchy and international unity were to remain; the center of unity was to keep its place. But the central power itself, its character, its origin, and its authority---all this was to be renewed.

The Romans themselves had a vague presentiment of this mysterious transformation. While the ordinary name of Rome was the Greek word for ``might,'' the citizens of the Eternal City believed that they discovered the true meaning of her name by reading it backward in Semitic fashion: \emph{amor}, and the ancient legend revived by Virgil connected the Roman people and the dynasty of Caesar in particular with the Mother of Love, and through her with the supreme God.
\end{quotationx}



\flrightit{Posted on 2011-01-08 by Cologero}

